% GENERATED FILE. Edit canonical lessons, then run npm run generate:books.
% canonical-source-hash: fnv1a64:8033b1e7599b34e4
% canonical-lessons: MR-C504-sutar, MR-C504-simpi, MR-C504-gavandi, MR-C504-nhavi, MR-C504-thokne

\chapter{Trades of the Town}
\label{ch:mr-trades-of-the-town}
\hlchaptermodality{504}

\begin{chapteropening}
\textbf{By the end of this chapter:} \emph{I can name some trades of the town and say what they do.}

Complete the last lesson of chapter 504: I can name some trades of the town and say what they do.
\end{chapteropening}

\section[\texorpdfstring{\mr{सुतार} (sutār)}{सुतार (sutār)}]{\mr{सुतार} (sutār) — a carpenter}

\label{lesson:MR-C504-sutar}



\begin{quote}
Before the new one: say the Marathi for dance, then the Marathi for a stage.
\end{quote}

\subsection*{You'll want to know: \mr{सुतार}}
\textbf{\mr{सुतार}} — \emph{sutār} — \textquotedblleft{}a carpenter\textquotedblright{}.

\begin{quote}
\mr{सुताराने} \mr{नवीन} \mr{दार} \mr{बनवलं}.
\end{quote}

\begin{quote}
\emph{sutārāne navīn dār banavlaṁ.}
\end{quote}

\begin{quote}
The carpenter made a new door.
\end{quote}

Masculine.

In the past the oblique \mr{सुतारा}- takes -\mr{ने}: \mr{सुताराने} \mr{दार} \mr{बनवलं}.

\begin{cousinweb}
\textbf{\mr{सुतार}} (\emph{sutār}) is a tadbhava, inherited from Sanskrit \emph{sūtradhāra}, \textquotedblleft{}carpenter\textquotedblright{} (CDIAL 13563), from \emph{sūtra}, \textquotedblleft{}a measuring line\textquotedblright{}, and \emph{dhāra}, \textquotedblleft{}holding\textquotedblright{}.
\end{cousinweb}

\begin{grammarlens}[title={What you've built}]
One more word for describing what you lived and felt.
\end{grammarlens}

\subsection*{Your turn}
\begin{itemize}
\raggedright
  \item \emph{Say it:} \emph{sutār}
  \item \emph{Say it:} \emph{sutār}, once more
  \item \emph{Say it:} \emph{sutārāne navīn dār banavlaṁ.} — The carpenter made a new door.
  \item \emph{Recall:} say the Marathi for dance, then the Marathi for a stage, then say \emph{sutār} again
\end{itemize}

\begin{tcolorbox}[breakable,colback=teal!4,colframe=teal!35!black,title={Before you move on}]
What does \mr{सुतार} mean? (\textquotedblleft{}A carpenter\textquotedblright{}.) Say it once more.
\end{tcolorbox}
\section[\texorpdfstring{\mr{शिंपी} (śimpī)}{शिंपी (śimpī)}]{\mr{शिंपी} (śimpī) — a tailor}

\label{lesson:MR-C504-simpi}



\begin{quote}
Before the new one: say the Marathi for a stage, then the Marathi for a carpenter.
\end{quote}

\subsection*{You'll want to know: \mr{शिंपी}}
\textbf{\mr{शिंपी}} — \emph{śimpī} — \textquotedblleft{}a tailor\textquotedblright{}.

\begin{quote}
\mr{शिंप्याने} \mr{माझा} \mr{सदरा} \mr{शिवला}.
\end{quote}

\begin{quote}
\emph{śimpyāne mājhā sadrā śivlā.}
\end{quote}

\begin{quote}
The tailor stitched my shirt.
\end{quote}

Masculine.

Before -\mr{ने} the -\mr{ई} of \mr{शिंपी} becomes -\mr{या}: \mr{शिंप्याने}.

\begin{cousinweb}
\textbf{\mr{शिंपी}} (\emph{śimpī}) is a tadbhava: CDIAL lists Marathi \emph{śĩpī}, the tailors' caste, under Sanskrit \emph{śilpin}, \textquotedblleft{}craftsman\textquotedblright{} (CDIAL 12471).
\end{cousinweb}

\begin{grammarlens}[title={What you've built}]
One more word for describing what you lived and felt.
\end{grammarlens}

\subsection*{Your turn}
\begin{itemize}
\raggedright
  \item \emph{Say it:} \emph{śimpī}
  \item \emph{Say it:} \emph{śimpī}, once more
  \item \emph{Say it:} \emph{śimpyāne mājhā sadrā śivlā.} — The tailor stitched my shirt.
  \item \emph{Recall:} say the Marathi for a stage, then the Marathi for a carpenter, then say \emph{śimpī} again
\end{itemize}

\begin{tcolorbox}[breakable,colback=teal!4,colframe=teal!35!black,title={Before you move on}]
What does \mr{शिंपी} mean? (\textquotedblleft{}A tailor\textquotedblright{}.) Say it once more.
\end{tcolorbox}
\section[\texorpdfstring{\mr{गवंडी} (gavaṇḍī)}{गवंडी (gavaṇḍī)}]{\mr{गवंडी} (gavaṇḍī) — a mason}

\label{lesson:MR-C504-gavandi}



\begin{quote}
Before the new one: say the Marathi for a carpenter, then the Marathi for a tailor.
\end{quote}

\subsection*{You'll want to know: \mr{गवंडी}}
\textbf{\mr{गवंडी}} — \emph{gavaṇḍī} — \textquotedblleft{}a mason\textquotedblright{}.

\begin{quote}
\mr{गवंड्याने} \mr{विटांची} \mr{भिंत} \mr{बांधली}.
\end{quote}

\begin{quote}
\emph{gavaṇḍyāne viṭāñcī bhint bāndhlī.}
\end{quote}

\begin{quote}
The mason built a brick wall.
\end{quote}

Masculine.

A builder in brick and stone; the oblique is \mr{गवंड्या}-, as in \mr{गवंड्याने}.

\begin{cousinweb}
\textbf{\mr{गवंडी}} (\emph{gavaṇḍī}) is the everyday Marathi word for a builder in brick and stone; its origin is not settled.
\end{cousinweb}

\begin{grammarlens}[title={What you've built}]
One more word for describing what you lived and felt.
\end{grammarlens}

\subsection*{Your turn}
\begin{itemize}
\raggedright
  \item \emph{Say it:} \emph{gavaṇḍī}
  \item \emph{Say it:} \emph{gavaṇḍī}, once more
  \item \emph{Say it:} \emph{gavaṇḍyāne viṭāñcī bhint bāndhlī.} — The mason built a brick wall.
  \item \emph{Recall:} say the Marathi for a carpenter, then the Marathi for a tailor, then say \emph{gavaṇḍī} again
\end{itemize}

\begin{tcolorbox}[breakable,colback=teal!4,colframe=teal!35!black,title={Before you move on}]
What does \mr{गवंडी} mean? (\textquotedblleft{}A mason\textquotedblright{}.) Say it once more.
\end{tcolorbox}
\section[\texorpdfstring{\mr{न्हावी} (nhāvī)}{न्हावी (nhāvī)}]{\mr{न्हावी} (nhāvī) — a barber}

\label{lesson:MR-C504-nhavi}



\begin{quote}
Before the new one: say the Marathi for a tailor, then the Marathi for a mason.
\end{quote}

\subsection*{You'll want to know: \mr{न्हावी}}
\textbf{\mr{न्हावी}} — \emph{nhāvī} — \textquotedblleft{}a barber\textquotedblright{}.

\begin{quote}
\mr{रविवारी} \mr{न्हाव्याच्या} \mr{दुकानात} \mr{मोठी} \mr{रांग} \mr{होती}.
\end{quote}

\begin{quote}
\emph{ravivārī nhāvyācyā dukānāt moṭhī rāṅg hotī.}
\end{quote}

\begin{quote}
On Sunday there was a long queue at the barber's shop.
\end{quote}

Masculine.

\begin{cousinweb}
\textbf{\mr{न्हावी}} (\emph{nhāvī}) is a tadbhava: CDIAL sets Marathi \emph{nhāvī} beside Sanskrit \emph{nāpita}, \textquotedblleft{}barber\textquotedblright{}, under \emph{*snāpitṛ}, \textquotedblleft{}one who bathes others\textquotedblright{} (CDIAL 13792).
\end{cousinweb}

\begin{grammarlens}[title={What you've built}]
One more word for describing what you lived and felt.
\end{grammarlens}

\subsection*{Your turn}
\begin{itemize}
\raggedright
  \item \emph{Say it:} \emph{nhāvī}
  \item \emph{Say it:} \emph{nhāvī}, once more
  \item \emph{Say it:} \emph{ravivārī nhāvyācyā dukānāt moṭhī rāṅg hotī.} — On Sunday there was a long queue at the barber's shop.
  \item \emph{Recall:} say the Marathi for a tailor, then the Marathi for a mason, then say \emph{nhāvī} again
\end{itemize}

\begin{tcolorbox}[breakable,colback=teal!4,colframe=teal!35!black,title={Before you move on}]
What does \mr{न्हावी} mean? (\textquotedblleft{}A barber\textquotedblright{}.) Say it once more.
\end{tcolorbox}
\section[\texorpdfstring{\mr{ठोकणे} (ṭhokṇe)}{ठोकणे (ṭhokṇe)}]{\mr{ठोकणे} (ṭhokṇe) — to hammer in; to knock}

\label{lesson:MR-C504-thokne}



\begin{quote}
Before the new one: say the Marathi for a mason, then the Marathi for a barber.
\end{quote}

\subsection*{You'll want to know: \mr{ठोकणे}}
\textbf{\mr{ठोकणे}} — \emph{ṭhokṇe} — \textquotedblleft{}to hammer in; to knock\textquotedblright{}.

\begin{quote}
\mr{सुताराने} \mr{भिंतीत} \mr{एक} \mr{खिळा} \mr{ठोकला}.
\end{quote}

\begin{quote}
\emph{sutārāne bhintīt ek khiḷā ṭhoklā.}
\end{quote}

\begin{quote}
The carpenter hammered a nail into the wall.
\end{quote}

Past, with -\mr{ने} (he …): \mr{त्याने} \mr{खिळा} \mr{ठोकला}, \emph{tyāne khiḷā ṭhoklā}.

The place takes -\mr{त}: \mr{भिंतीत} \mr{खिळा} \mr{ठोकणे}, to drive a nail into a wall.

\begin{cousinweb}
\textbf{\mr{ठोकणे}} (\emph{ṭhokṇe}) is inherited: CDIAL lists Marathi \emph{ṭhokṇẽ}, \textquotedblleft{}to strike\textquotedblright{}, under \emph{*ṭhōkk}, \textquotedblleft{}knock\textquotedblright{} (CDIAL 5513), a sound word shared with Hindi \emph{ṭhoknā}.
\end{cousinweb}

\begin{grammarlens}[title={What you've built}]
One more word for describing what you lived and felt.
\end{grammarlens}

\subsection*{Your turn}
\begin{itemize}
\raggedright
  \item \emph{Say it:} \emph{ṭhokṇe}
  \item \emph{Say it:} \emph{ṭhokṇe}, once more
  \item \emph{Say it:} \emph{sutārāne bhintīt ek khiḷā ṭhoklā.} — The carpenter hammered a nail into the wall.
  \item \emph{Recall:} say the Marathi for a mason, then the Marathi for a barber, then say \emph{ṭhokṇe} again
\end{itemize}

\begin{tcolorbox}[breakable,colback=teal!4,colframe=teal!35!black,title={Before you move on}]
What does \mr{ठोकणे} mean? (\textquotedblleft{}To hammer in; to knock\textquotedblright{}.) Say it once more.
\end{tcolorbox}
